\chapter{The Murch}

The murch painted by Tibor McMasters spread slowly across the world. At first it circulated as rumor, then as reproduction, then as destination. In time it came to be judged equal to the works of the great Italian masters, most of whose originals no longer existed and were known only through imperfect prints. In this sense, Tibor’s work belonged naturally among them: an image surviving a ruined age.

Seventeen years after Tibor’s death, the hierarchy of the Servants of Wrath issued a final and binding declaration of authenticity. The image was proclaimed to be, beyond question, the true visage of the God of Wrath—Carleton Lufteufel. All debate ended there. To dispute the identification was henceforth illegal. The penalties were explicit: emasculation for men, removal of one ear for women. These measures were instituted to secure reverence in an irreverent world, faith in a society that had misplaced it, belief in a culture that had learned—too late—that much of what it once believed had been lies.

At the time of his death, Tibor subsisted on a modest annual pension from the Church, supplemented by guaranteed upkeep of his cart and fodder for two cows. Because of the excellence of his work he had been granted two animals instead of one. As he traveled, people recognized him and called out to him. Tourists requested autographs, which he gave with effort and patience. Children shouted his name, but did not jeer. He was liked. Though he grew eccentric and irritable with age, he was regarded as an asset to the community. This esteem endured despite the fact that, after completing the true portrait of the God of Wrath, Tibor McMasters never again produced a work of significance.

It was rumored that among his personal effects were fragments of a diary—private notes written only for himself—in which, late in life, he recorded certain doubts regarding the authenticity of his own great murch. No such documents were ever produced. If they existed at all, the Servants of Wrath, who took possession of his papers after his death, either sealed them behind locked metal doors or, more probably, destroyed them.

Tibor’s last two cows were slaughtered, preserved, and mounted—one on each side of the murch—to stare solemnly, glass-eyed, at the pilgrims and tourists who came to pay homage to the painting. Tibor McMasters himself was eventually canonized by the Church. His grave site is unknown. Several cities claim it, proudly.

And so he entered history as the man who had seen God, and shown Him—
and as the one who, having done so, could paint nothing more.

